\subsection{Detector Response}\label{sec:RECONSTRUCTION}
The signal-processed waveforms are then passed forward to the reconstruction routines. A peak finding algorithm fits Gaussian distributions to the signal waveforms, where each successful fit is formed into an object known as a \gls{tpc} \textit{hit}. The area under all the \gls{tpc} hits is proportional to the energy deposited by the charge particles as they propagate through the liquid argon volume.  For the horizontal drift processing chain, a disambiguation routine is then applied to the induction plane hits to deduce which drift volume they originated from. This routine combines spatial and timing information of all hits in all planes simultaneously to maximize performance. The \gls{tpc} hits are then passed forward to the higher-level reconstruction routines.

The propagation of scintillation photons in liquid argon is not performed on a per-event basis, as it would be too computationally demanding giving the high \gls{lar} light-yield, which is about 24000 photons per MeV of energy deposited for a \gls{mip}. Instead, the propagation considering all possible effects (Rayleigh scattering, absorption by impurities, reflection and absorption on surfaces) is performed systematically in order to obtain fits used to predict the number of photons that should land on the sensitive surface of the photon detectors and their arrival time at run time. This {\it fast optical simulation} model is described in Ref.~\cite{semiana} and provides a proper dependency of the detection light-yield on the distance and offset angle between the origin of the photon and the center of a \gls{pd}. 

Once the number of arriving photons as a function of time is obtained per \gls{pd} channel, the digitization of the signal is performed. The generation of optical waveforms takes into account the expected \gls{pde} of the \gls{xa}s, the single \gls{pe} signal response, \gls{sipm} crosstalk, dark count rate, and electronic noise. Those waveforms go through a peak finding algorithm. For the \gls{vd} case, this is simply done by looking for peaks above a threshold (usually corresponding to 1.5 \gls{pe}s) and assigning the number of \gls{pe}s as the integral of the peak divided by the area corresponding to a single \gls{pe}. However, the \gls{hd} case presents additional complexity, as the single \gls{pe} shape shows a sizeable undershoot which requires a deconvolution to be performed before obtaining the number of \gls{pe}s in a given peak.  

Both reconstruction methods developed for low-energy neutrino analyses in \gls{dune} are presented in the following chapters. Raw \gls{tpc} hits, obtained after the aforementioned detector response stage, serve as the starting point for both workflows. These reconstruction objects result from the deconvolution, noise reduction and Gaussian fitting steps of the raw waveforms. Accounting for the signal's vertex of interaction within the detector volume, the \gls{tpc} hit charge can be corrected, as the free electron cloud diffuses and decreases with a characteristic lifetime while propagating through the drift coordinate. This results in a lower imprint onto the \gls{tpc}'s collection plane depending on the signal projection distance.

Beginning at the hit level, we apply calibration steps to calculate an energy that accounts for detector response and physical processes. These include ionisation, diffusion, electron lifetime, and later, clustering losses:

\begin{itemize}
    \item \textbf{Recombination factor}. A constant factor (0.63) is applied a priori to the integral fit of the readout waveform to account for the amount of energy that is immediately reabsorbed by the ionisation electrons by argon atoms, producing scintillation light and reducing the induced and collected charge. This value is related to the electric field strength, causing ionizing electrons to drift towards the anode.% The value used here is a common value used in other studies.
    \item \textbf{Plane Calibration}. A factor is applied to each readout, U, V and W, to correct for the difference in signal current between the induced and collected charge on the final readout waveform.
    \item \textbf{Electron lifetime}. A correction is calculated from the exponential decay of the collected charge divided by the true electron energy as a function of the adjusted detector time and drift distance. The simulated electron lifetime is $\SI{10}{\milli\second}$, lower than the one measured in ProtoDUNE-HD~\cite{DUNE:2020cqd}.%  Due to collection and threshold effects, the fitted value turns out to be lower than the true input of the simulation. This depends on the length of the drift and background contamination.
\end{itemize}